\documentclass[11pt,a4paper]{article}
\usepackage[margin=2.5cm,headheight=15pt]{geometry}
\usepackage{amsmath,amssymb,amsthm}
\usepackage{mathtools}
\usepackage{microtype}
\usepackage{hyperref}
\usepackage{cleveref}
\usepackage{enumitem}
\usepackage{mdframed}
\usepackage{xcolor}
\usepackage{titlesec}
\usepackage{booktabs}
\usepackage{fancyhdr}

\definecolor{navyblue}{RGB}{26,58,92}
\definecolor{midblue}{RGB}{44,95,138}
\definecolor{lightblue}{RGB}{232,240,248}
\definecolor{proofgray}{RGB}{242,242,242}

\hypersetup{colorlinks=true,linkcolor=midblue,citecolor=midblue,urlcolor=midblue}

\pagestyle{fancy}\fancyhf{}
\fancyhead[L]{\small\color{navyblue}\textit{Non-Compact Configuration-Based Formulation for the U-SSP}}
\fancyhead[R]{\small\color{navyblue}\thepage}
\renewcommand{\headrulewidth}{0.4pt}

\titleformat{\section}{\large\bfseries\color{navyblue}}{}{0em}{}[\color{navyblue}\titlerule]
\titleformat{\subsection}{\normalsize\bfseries\color{midblue}}{}{0em}{}
\titleformat{\subsubsection}{\normalsize\itshape\color{midblue}}{}{0em}{}

\newmdenv[backgroundcolor=lightblue,linecolor=navyblue,linewidth=1.2pt,
  leftline=true,rightline=false,topline=false,bottomline=false,
  innerleftmargin=12pt,innerrightmargin=8pt,innertopmargin=8pt,
  innerbottommargin=8pt,skipabove=10pt,skipbelow=6pt]{thmbox}

\newmdenv[backgroundcolor=proofgray,linecolor=midblue,linewidth=0.8pt,
  leftline=true,rightline=false,topline=false,bottomline=false,
  innerleftmargin=10pt,innerrightmargin=8pt,innertopmargin=6pt,
  innerbottommargin=6pt,skipabove=6pt,skipbelow=6pt]{proofbox}

\theoremstyle{plain}
\newtheorem{theorem}{Theorem}[section]
\newtheorem{lemma}[theorem]{Lemma}
\newtheorem{proposition}[theorem]{Proposition}
\newtheorem{corollary}[theorem]{Corollary}
\theoremstyle{definition}
\newtheorem{definition}[theorem]{Definition}
\newtheorem{remark}[theorem]{Remark}
\newtheorem{example}[theorem]{Example}

\setlist[itemize]{noitemsep,topsep=4pt}

% ─────────────────────────────────────────────────────────────────────────────
\begin{document}
% ─────────────────────────────────────────────────────────────────────────────

\begin{center}
  {\color{navyblue}\rule{\textwidth}{2pt}}\\[0.4cm]
  {\LARGE\bfseries\color{navyblue}The Job Sequencing and Tool Switching Problem}\\[0.3cm]
  {\large\itshape\color{midblue}Part II: The Non-Compact Configuration-Based Formulation}\\[0.2cm]
  {\color{navyblue}\rule{\textwidth}{2pt}}
\end{center}

\begin{mdframed}[backgroundcolor=lightblue,linecolor=navyblue,linewidth=1pt,
  innerleftmargin=14pt,innerrightmargin=14pt,innertopmargin=8pt,innerbottommargin=8pt]
\small\textbf{Scope.}
This document presents the \emph{non-compact, configuration-based} ILP formulation
for the Uniform SSP, proposed by Colares and Wagler \cite{Colares2026Exact}.
It includes: the configuration space and job clusters, the transition metric with
a metric-space proof, the full ILP with PCTSP-based subtour elimination (which is
tighter than GSECs), and the column generation framework with the H-SUKP pricing
subproblem. This is the primary modelling framework for all subsequent analysis.

\smallskip
\textbf{Key caveat}: the PCTSP-based formulation requires \emph{both} column
generation (exponential variables) and row generation (exponential constraints),
creating a ``double generation'' problem. Part~IV (\texttt{04\_mtz\_formulation.tex})
resolves this by replacing the exponential constraints with polynomial MTZ constraints,
leaving only column generation.
\end{mdframed}

\vspace{0.5cm}
\tableofcontents
\newpage

% ─────────────────────────────────────────────────────────────────────────────
\section{Motivation: Escaping Compact-Model Failure}
\label{sec:motivation}
% ─────────────────────────────────────────────────────────────────────────────

As established in \texttt{01\_foundations.tex}, all compact ILP formulations
for the SSP suffer from two fundamental problems:
(i) the LP relaxation collapses to a lower bound of 0, providing no pruning
information, and
(ii) the 0-block symmetry creates $2^k$ equivalent integer solutions that
exhaust the branch-and-bound tree.

The non-compact formulation introduced here replaces positional variables with
explicit \emph{configuration} variables — one per valid magazine state. This
eliminates the symmetry problem and gives a much tighter LP relaxation. However,
it introduces a new computational challenge that must be understood upfront.

\begin{mdframed}[backgroundcolor=lightblue,linecolor=navyblue,linewidth=1pt,
  innerleftmargin=14pt,innerrightmargin=14pt,innertopmargin=8pt,innerbottommargin=8pt]
\textbf{The double-generation problem.}
The non-compact formulation has \emph{two} sources of exponential size
simultaneously:
\begin{itemize}
\item \textbf{Exponential columns} (configurations): $|\mathcal{C}| = \binom{m}{b}$
  variables. Handled by \textbf{column generation}: add improving configurations
  one at a time via a pricing oracle (H-SUKP, Section~\ref{sec:sukp}).
\item \textbf{Exponential rows} (subtour constraints): the PCTSP-based subtour
  elimination constraints number $O(2^{|\mathcal{C}|})$. Handled by
  \textbf{row generation}: add violated constraints on-the-fly via separation.
\end{itemize}
Running column generation and row generation \emph{simultaneously} is called
\emph{double generation}. It is technically complex: each LP re-solve after
adding a new column potentially violates previously satisfied subtour cuts,
requiring a full re-separation pass. This interaction makes convergence slow and
implementation fragile.

\medskip
\textbf{The remedy} (Part~IV, \texttt{04\_mtz\_formulation.tex}):
replace the exponential subtour constraints with a \emph{polynomial} set of
Miller--Tucker--Zemlin (MTZ) constraints. This eliminates row generation entirely,
leaving only column generation — a much more tractable problem.
The trade-off is a slightly weaker LP bound, but in practice this is far
outweighed by the reduction in implementation complexity and solver iterations.
\end{mdframed}

\medskip
The non-compact approach takes a different viewpoint from compact models: instead
of asking ``which tool is at position $k$?'', it asks ``which complete magazine
state (configuration) is used to process a group of jobs?'' By working directly
over the set of all valid magazine states, the model eliminates positional
symmetry entirely and achieves a far tighter LP relaxation.

% ─────────────────────────────────────────────────────────────────────────────
\section{Configuration Space and Job Clusters}
\label{sec:config_space}
% ─────────────────────────────────────────────────────────────────────────────

We use the same notation as \texttt{01\_foundations.tex}:
$J=\{1,\ldots,n\}$ jobs, $T=\{1,\ldots,m\}$ tools, magazine capacity $b$,
requirement sets $T_j \subseteq T$.

\begin{definition}[Configuration space $\mathcal{C}$]
The \emph{configuration space} is the set of all maximal tool subsets:
\[
  \mathcal{C} = \bigl\{ C \subseteq T \;\big|\; |C| = b \bigr\}.
\]
Its cardinality is $|\mathcal{C}| = \binom{m}{b}$.
\end{definition}

By Lemma~1 of \texttt{01\_foundations.tex}, restricting to maximal configurations
is without loss of generality. We augment $\mathcal{C}$ with a dummy configuration
$C_0 = \emptyset$ (or any fixed initial state) representing the machine's magazine
before any job is processed.

\begin{definition}[Job cluster $H_j$]
For each job $j \in J$, its \emph{cluster} is the set of configurations that
can process it:
\[
  H_j = \bigl\{ C \in \mathcal{C} \;\big|\; T_j \subseteq C \bigr\}.
\]
Since $|T_j| \le b$, every cluster is non-empty.
\end{definition}

\noindent\textbf{Key point.}
Unlike in a standard GTSP, clusters can \emph{overlap}: a single configuration
$C$ may satisfy multiple jobs simultaneously when $T_i \cup T_j \subseteq C$.
This overlap is the central structural property of the SSP and is handled by
the GTSP equivalence in \texttt{03\_gtsp\_equivalence.tex}.

% ─────────────────────────────────────────────────────────────────────────────
\section{The Transition Metric}
\label{sec:metric}
% ─────────────────────────────────────────────────────────────────────────────

\subsection{Definition}

When the machine transitions from a magazine state $C_u$ to $C_v$, the tools
in $C_u \setminus C_v$ are evicted and the tools in $C_v \setminus C_u$ must
be inserted. Since both configurations are maximal ($|C_u|=|C_v|=b$), the number
of evictions equals the number of insertions, and the switching cost is:
\[
  d(C_u, C_v) = |C_v \setminus C_u| = b - |C_u \cap C_v|.
\]

\subsection{Metric Space Validity}

\begin{thmbox}
\begin{theorem}[Triangle inequality for $d$ \cite{Colares2026Exact}]
\label{thm:triangle}
For any three configurations $C_u, C_w, C_v \in \mathcal{C}$:
\[
  d(C_u, C_v) \le d(C_u, C_w) + d(C_w, C_v).
\]
\end{theorem}
\end{thmbox}

\begin{proofbox}
\begin{proof}
Using $d(A,B) = |A \setminus B|$, we show $|C_u \setminus C_v| \le
|C_u \setminus C_w| + |C_w \setminus C_v|$.

Pick any tool $t \in C_u \setminus C_v$, so $t \in C_u$ and $t \notin C_v$.
\begin{itemize}
\item If $t \notin C_w$: then $t \in C_u \setminus C_w$.
\item If $t \in C_w$: then since $t \notin C_v$, we have $t \in C_w \setminus C_v$.
\end{itemize}
In both cases $t \in (C_u \setminus C_w) \cup (C_w \setminus C_v)$, so
$C_u \setminus C_v \subseteq (C_u \setminus C_w) \cup (C_w \setminus C_v)$.
Taking cardinalities and applying subadditivity of set union:
\[
  |C_u \setminus C_v| \le |(C_u \setminus C_w) \cup (C_w \setminus C_v)|
  \le |C_u \setminus C_w| + |C_w \setminus C_v|. \qquad\qed
\]
\end{proof}
\end{proofbox}

\begin{remark}
The metric is also symmetric ($d(C_u,C_v)=d(C_v,C_u)$) and non-negative,
so $(\mathcal{C}, d)$ is a valid metric space. Symmetry follows from
$|A \setminus B| = |B \setminus A|$ when $|A|=|B|=b$.
The maximum distance is $b$ (completely disjoint configurations).
\end{remark}

% ─────────────────────────────────────────────────────────────────────────────
\section{The Non-Compact ILP}
\label{sec:nc_ilp}
% ─────────────────────────────────────────────────────────────────────────────

\subsection{Decision Variables}

\begin{itemize}
\item $y(v) \in \{0,1\}$: 1 if configuration $v \in \mathcal{C}$ is \emph{activated}
      (used at least once in the sequence).
\item $x(u,v) \in \{0,1\}$: 1 if the machine transitions \emph{directly} from
      configuration $u$ to configuration $v$.
\end{itemize}

We include $C_0$ (dummy start/end node) with $y(C_0)=1$ fixed and $d(C_0,v)$
equal to the cost of loading configuration $v$ from an empty magazine.

\subsection{Objective and Base Constraints}

\begin{align}
\min \quad & \sum_{u \in \mathcal{C}} \sum_{v \in \mathcal{C} \setminus \{u\}}
             d(u,v)\, x(u,v) \label{eq:nc_obj}\\[4pt]
\text{s.t.}\quad
& \sum_{v \in H_j} y(v) \ge 1 \quad \forall j \in J
  \label{eq:nc_cover}\\
& \sum_{u \ne v} x(u,v) = y(v) \quad \forall v \in \mathcal{C}
  \label{eq:nc_flowin}\\
& \sum_{v \ne u} x(u,v) = y(u) \quad \forall u \in \mathcal{C}
  \label{eq:nc_flowout}\\
& y(C_0) = 1 \label{eq:nc_depot}\\
& x(u,v) \in \{0,1\},\; y(v) \in \{0,1\} \quad \forall u,v \in \mathcal{C}
  \notag
\end{align}

\noindent\textbf{Interpretation.}
Constraint~\eqref{eq:nc_cover} ensures every job is processed by at least one
activated configuration. Constraints~\eqref{eq:nc_flowin}--\eqref{eq:nc_flowout}
enforce flow conservation: each activated node has exactly one incoming and one
outgoing arc, forming a collection of paths/cycles. The depot constraint~\eqref{eq:nc_depot}
anchors the sequence to start and end at $C_0$.

\subsection{Subtour Elimination: PCTSP-Based SECs}

Constraints~\eqref{eq:nc_flowin}--\eqref{eq:nc_flowout} alone allow disconnected
subtours among activated configurations. We need constraints that force all
activated configurations to lie on a \emph{single} path through the depot.

The standard approach uses Generalized Subtour Elimination Constraints (GSECs).
However, GSECs are only valid when $y(v) \ge 1$ for all nodes in the subset,
which does not hold here since $y(v) \in \{0,1\}$. A tighter and correct
formulation uses the Prize-Collecting TSP (PCTSP) subtour elimination constraints
\cite{Balas1989}:
\[
  \sum_{u \in S}\sum_{v \in S \setminus \{u\}} x(u,v)
  \;\le\; \sum_{u \in S \setminus \{k\}} y(u)
  \quad \forall S \subset \mathcal{C},\; 2 \le |S| \le |\mathcal{C}|-1,\;
  C_0 \notin S,\; \forall k \in S.
  \label{eq:nc_pctsp}
\]

\begin{remark}[Why PCTSP over GSEC]
The GSEC form $\sum_{u,v \in S} x(u,v) \le |S|-1$ assumes all $y(u)=1$ within
$S$. When some $y(u)=0$ (unactivated), the bound is too loose and subtours of
inactive nodes can appear in LP solutions. The PCTSP form $\sum x(u,v) \le
\sum_{u \ne k} y(u)$ correctly adapts to the actual number of \emph{activated}
nodes, eliminating spurious subtours even in the fractional LP relaxation.
See \cite{Balas1989} for the original derivation.
\end{remark}

\subsection{Complete Non-Compact ILP}

Combining all components:
\begin{align}
\min \quad & \sum_{u,v} d(u,v)\, x(u,v) \notag\\
\text{s.t.}\quad
& \sum_{v \in H_j} y(v) \ge 1 && \forall j \in J \notag\\
& \sum_{u \ne v} x(u,v) = y(v) && \forall v \in \mathcal{C} \notag\\
& \sum_{v \ne u} x(u,v) = y(u) && \forall u \in \mathcal{C} \notag\\
& \sum_{u \in S}\sum_{v \in S \setminus \{u\}} x(u,v)
  \le \sum_{u \in S \setminus \{k\}} y(u)
  && \forall S \subset \mathcal{C},\, C_0 \notin S,\, \forall k \in S \notag\\
& y(C_0) = 1 \notag\\
& x(u,v),\, y(v) \in \{0,1\} && \forall u,v \in \mathcal{C} \notag
\end{align}

\noindent Since $|\mathcal{C}| = \binom{m}{b}$ is exponential, the full model
cannot be instantiated explicitly. This motivates the column generation approach
in the next section.

% ─────────────────────────────────────────────────────────────────────────────
\section{Column Generation Framework}
\label{sec:cg}
% ─────────────────────────────────────────────────────────────────────────────

\subsection{Restricted Master Problem (RMP)}

Let $V \subset \mathcal{C}$ be a small, dynamically growing subset of
configurations (initially seeded by heuristic solutions or greedy construction).
The RMP solves the LP relaxation of the non-compact ILP restricted to $V$:
\begin{align}
\min \quad & \sum_{i,j \in V} d(i,j)\, x_{ij} \label{eq:rmp_obj}\\
\text{s.t.}\quad
& \sum_{i \in V \cap H_r} y_i \ge 1 && \forall r \in J
  \quad (\text{dual: } \pi_r \ge 0) \label{eq:rmp_cover}\\
& \sum_{j \ne i} x_{ij} = y_i && \forall i \in V
  \quad (\text{dual: } \alpha_i \in \mathbb{R}) \label{eq:rmp_flowin}\\
& \sum_{j \ne i} x_{ji} = y_i && \forall i \in V
  \quad (\text{dual: } \beta_i \in \mathbb{R}) \label{eq:rmp_flowout}\\
& \sum_{i \in S}\sum_{j \in S \setminus\{i\}} x_{ij}
  \le \sum_{i \in S \setminus\{k\}} y_i
  && \forall S \subset V,\, C_0 \notin S,\, k \in S
  \quad (\text{dual: } \mu_{S,k} \le 0) \label{eq:rmp_sec}\\
& y_{C_0} = 1 \notag\\
& x_{ij} \ge 0,\; y_i \ge 0 && \forall i,j \in V \notag
\end{align}

At each iteration, solving the RMP yields LP optimal value $Z^*_{RMP}$ and dual
variables $(\pi, \alpha, \beta, \mu)$.

\subsection{Pricing Subproblem: When to Add a Configuration}

A new configuration $c \notin V$ should be added if doing so can reduce the
RMP objective. The \emph{reduced cost} of introducing $c$ with incoming arc from
$u \in V$ and outgoing arc to $w \in V$ is:
\[
  \bar{c}(u,c,w) = d(u,c) + d(c,w) - \alpha_u - \beta_w
  - \sum_{r \in J:\, T_r \subseteq c} \pi_r
  - \sum_{S,k:\, u,c \in H(S)} \mu_{S,k}
\]
(the GSEC dual terms vanish for a new node $c$ not yet in any identified subset
$S$). A configuration $c$ is \emph{improving} if $\min_{u,w} \bar{c}(u,c,w) < 0$.

Ignoring routing duals for the pricing step (standard practice), the dominant
term is the coverage gain. We seek:
\[
  \max_{c:\, |c|=b} \;\sum_{r \in J:\, T_r \subseteq c} \pi_r
\]
subject to $|c| \le b$ and the configuration feasibility condition
$T_r \subseteq c \iff v_r = 1$ (all tools of job $r$ are in $c$).

\subsection{The Pricing Subproblem as a Hypergraph SUKP}
\label{sec:sukp}

Let $z_t \in \{0,1\}$ indicate that tool $t$ is included in the new configuration,
and $v_r \in \{0,1\}$ indicate that job $r$ is \emph{covered} by the configuration
(meaning all its required tools are selected). The pricing problem is:

\begin{align}
\max \quad & \sum_{r \in J} \pi_r\, v_r \label{eq:ps_obj}\\
\text{s.t.}\quad
& v_r \le z_t && \forall r \in J,\, \forall t \in T_r \label{eq:ps_cover}\\
& \sum_{t \in T} z_t \le b \label{eq:ps_cap}\\
& z_t \in \{0,1\} \;\; \forall t \in T, \quad v_r \in \{0,1\} \;\; \forall r \in J
  \notag
\end{align}

Constraint~\eqref{eq:ps_cover} ensures that $v_r=1$ only if all tools in $T_r$
are selected. This is the \emph{Set Union Knapsack Problem (SUKP)} with hyperedges
(one hyperedge per job, connecting its required tools) and item capacity $b$:
it is \textbf{strongly NP-hard} \cite{Goldschmidt1994}. It can be solved exactly
by branch-and-bound or IP, with ML-based acceleration for large instances
(see \texttt{08\_research\_notes.md}).

\begin{remark}[SUKP vs.\ standard knapsack]
In a standard knapsack, each item contributes independently to value and weight.
In the SUKP, a job $r$ contributes reward $\pi_r$ only if \emph{all} its tools
are selected simultaneously — a non-separable, conjunctive structure. This is
why the pricing subproblem is strictly harder than a standard knapsack.
\end{remark}

\subsection{Column Generation Termination}

The LP relaxation is solved to optimality when no improving column can be found:
$\min_{c} \bar{c}(u,c,w) \ge 0$ for all potential insertion pairs $(u,w)$.
This certificate is provided by the exact SUKP solver. If a negative reduced
cost column is found, it is added to $V$ and the RMP is re-solved.
The process continues until optimality is certified.
Integer solutions are recovered by branch-and-price \cite{Barnhart1998}.

% ─────────────────────────────────────────────────────────────────────────────
\section{Why the Non-Compact Formulation is Better}
\label{sec:advantages}
% ─────────────────────────────────────────────────────────────────────────────

\begin{enumerate}
\item \textbf{No positional symmetry.} Each configuration is a unique physical
  machine state, not a positional variable. The $2^k$ 0-block symmetry of compact
  models disappears: the KTNS policy is implicit in the choice of which
  configurations appear on the path.
\item \textbf{Tight LP relaxation.} The coverage constraints over job clusters,
  combined with flow conservation, provide a much stronger lower bound than the
  zero bound of position-based models.
\item \textbf{Scalable via column generation.} We never enumerate all
  $\binom{m}{b}$ configurations; only those with improving reduced costs are added.
\item \textbf{Natural GTSP structure.} The non-compact ILP is isomorphic to a
  Generalised TSP over the configuration graph (Part~III), enabling the use of
  decades of TSP theory and solvers.
\item \textbf{The double-generation caveat.} The PCTSP-based formulation
  above requires both column generation and row generation simultaneously.
  Part~IV (\texttt{04\_mtz\_formulation.tex}) resolves this by replacing the
  exponential PCTSP constraints with $O(n^2)$ MTZ constraints, eliminating
  row generation entirely at the cost of a slightly weaker LP bound.
\end{enumerate}


\bibliographystyle{abbrvnat}
\bibliography{references}
\end{document}
